% =============================================================================
% ANEXOS (\appendix activado en main.tex)
% =============================================================================

\chapter{Visualizaciones complementarias}
\label{anexo:radares}

Este anexo recoge las visualizaciones complementarias del cuerpo
principal: las fichas \textit{scout-facing} por jugador para tres
casos paradigmáticos del torneo distintos al de Mbappé reportado en
la figura~\ref{fig:radar-mbappe}, los \textit{scatter} de la selección
de Francia sobre las dos combinaciones de índices del \CCV, la figura
\textit{event-study} de la verificación de robustez en
diferencias-en-diferencias y las cotas \textit{HonestDiD} de
sensibilidad a violaciones de paralelismo.

\section{Fichas \textit{scout-facing} adicionales}
\label{anexo:fichas}

Las figuras \ref{fig:radar-messi}, \ref{fig:radar-hakimi} y
\ref{fig:radar-brozovic} reportan la ficha equivalente para tres
jugadores con perfiles contextuales distintos al de Mbappé. La
lectura aplica el mismo criterio que la figura~\ref{fig:radar-mbappe}
del cuerpo principal: el percentil que muestra cada eje es la
posición relativa del jugador dentro de su \textit{bucket} posicional
sobre el corpus del torneo, no una medida de significatividad
estadística del efecto causal sobre el jugador. La probabilidad
posterior por jugador y celda, base de los \textit{tiers} de
significatividad, se reporta por separado en la tabla canónica del
\CCV.

\begin{figure}[h]
  \centering
  \includegraphics[width=0.95\textwidth]{radar_report_1531}
  \caption{Ficha \textit{scout-facing} de Lionel Messi.}
  \label{fig:radar-messi}
\end{figure}

\begin{figure}[h]
  \centering
  \includegraphics[width=0.95\textwidth]{radar_report_1681}
  \caption{Ficha \textit{scout-facing} de Achraf Hakimi.}
  \label{fig:radar-hakimi}
\end{figure}

\begin{figure}[h]
  \centering
  \includegraphics[width=0.95\textwidth]{radar_report_8129}
  \caption{Ficha \textit{scout-facing} de Marcelo Brozović.}
  \label{fig:radar-brozovic}
\end{figure}

\section{Vista por equipo: Francia}
\label{anexo:france-scatters}

Las figuras \ref{fig:scatter-france-remontador-cerrojo} y
\ref{fig:scatter-france-ataque} reproducen los dos \textit{scatter}
del cuerpo principal restringidos a los jugadores de la selección
campeona, con la nube del resto del torneo como contexto de fondo.
La lectura por equipo permite identificar la posición relativa de
cada jugador del bloque francés sobre el panorama agregado.

\begin{figure}[h]
  \centering
  \includegraphics[width=\textwidth]{scatter_france_remontador_cerrojo}
  \caption{Vista por equipo (Francia): Remontador frente a Cerrojo.}
  \label{fig:scatter-france-remontador-cerrojo}
\end{figure}

\begin{figure}[h]
  \centering
  \includegraphics[width=\textwidth]{scatter_france_ataque_marcar_encajar}
  \caption{Vista por equipo (Francia): ataque tras marcar frente a
    ataque tras encajar.}
  \label{fig:scatter-france-ataque}
\end{figure}

\section{Verificación causal por \textit{event-study}}
\label{anexo:event-study}

La figura~\ref{fig:event-study} reporta los coeficientes del estudio
de eventos por minuto relativo al \textit{shock}, estimados mediante
el estimador robusto a heterogeneidad de \textcite{SunAbraham2021}.
La ausencia de tendencia sistemática en los minutos pre-\textit{shock}
es consistente con el test $F$ de pre-tendencias agregadas
\parencite{Roth2022Pretest} reportado en el capítulo~\ref{ch:resultados}.

\begin{figure}[h]
  \centering
  \includegraphics[width=\textwidth]{event_study}
  \caption{\textit{Event-study} por minuto relativo al \textit{shock}.}
  \label{fig:event-study}
\end{figure}

\section{Sensibilidad \textit{HonestDiD} a violaciones de paralelismo}
\label{anexo:honestdid}

La figura~\ref{fig:honestdid} reporta las cotas de sensibilidad de
\textcite{RambachanRoth2023} sobre el ATE poblacional para tres
niveles de relajación de la asunción de tendencias paralelas. Los
intervalos contienen el cero en todos los casos, lo que es
consistente con la indistinguibilidad de cero reportada para el ATE
en el capítulo~\ref{ch:resultados}.

\begin{figure}[h]
  \centering
  \includegraphics[width=\textwidth]{fig_anexo_honestdid}
  \caption{Cotas \textit{HonestDiD} sobre el ATE poblacional.}
  \label{fig:honestdid}
\end{figure}


\chapter{Limitaciones metodológicas en detalle}
\label{anexo:limitaciones}

Este anexo amplía las cautelas declaradas en la
sección~\ref{sec:disc-no-demuestran} del capítulo~\ref{ch:discusion}
con tres familias de limitaciones que el lector debe considerar para
no sobre-interpretar los resultados.

\section{Implementación propia frente a librerías de referencia}
\label{anexo:lim-implementacion}

El canal ofensivo utiliza \textit{atomic}-VAEP a través de
\texttt{socceraction} 1.5.3 \parencite{Decroos2019VAEP}, la
implementación oficial mantenida por los autores originales. Los
demás \textit{building blocks} ---PPCF de
\textcite{Spearman2018OBSO}, C-OBSO de \textcite{Teranishi2022COBSO},
VDEP estricto de \textcite{Toda2022VDEP}, \textit{exPress} de
\textcite{Lee2025exPress} y la atribución \textit{frame-level} al
defensor más cercano--- son reimplementaciones propias derivadas
de la lectura directa de los \textit{papers} originales. La validación
es triple: (i) reproducción de las cifras agregadas reportadas en los
\textit{papers}; (ii) comparación con \textit{outputs} comerciales
(PFF FC, StatsBomb) en corpus solapados; (iii) consistencia interna
entre canales sobre los mismos jugadores. El nivel es razonable para
un trabajo de máster, pero las reimplementaciones no tienen la
madurez ni la adopción comunitaria de una librería oficial; cualquier
hallazgo numérico que dependa de manera crítica de la implementación
exacta debería contrastarse con una segunda implementación cuando
esté disponible.

\section{Cobertura limitada al Mundial Qatar 2022}
\label{anexo:lim-cobertura}

El \CCV se aplica al Mundial Qatar 2022 por ser el único corpus
público que combina simultáneamente \textit{eventing} sincronizado,
tracking a 25\,Hz, cobertura completa de fase eliminatoria,
\textit{grades} humanos por evento y rosters oficiales. Las
propiedades individuales identificadas corresponden a este torneo
concreto, no son una predicción para otros contextos. Los 172
\textit{shocks} sobre 64 partidos sitúan el tamaño de muestra
efectivo en el extremo bajo de lo que la identificación causal a
nivel individual admite.

\section{Sub-corpus de \textit{near-miss} reducido}
\label{anexo:lim-nearmiss}

El cuasi-experimento de \textit{near-miss} depende de 70 disparos
efectivos sobre las cinco tipologías pre-registradas. Es suficiente
para triangular con el DiD a nivel agregado, pero no para
identificar efectos individuales con poder estadístico. La
estrategia conservadora es priorizar el DiD como brazo principal;
la triangulación rigurosa por jugador queda como línea futura sobre
corpus mayores.

\section{Qué \textit{no} captura el \CCV}
\label{anexo:lim-no-captura}

Por construcción, el marco deja fuera cinco fenómenos relevantes
para la lectura causal completa.

\paragraph{Picos concentrados de uno o dos minutos.}
El CATE modela respuesta sostenida en la ventana $\pm 10$\,min. Un
jugador con una jugada decisiva concentrada en uno o dos minutos
cuyo delta promedio en la ventana es modesto queda subestimado. El
caso paradigmático es Mbappé en los minutos 80--82 de la final, con
el doblete en 97 segundos: el marco lo etiqueta como
\textit{pressure-clutch} significativo, pero la magnitud del CATE
infraponderada respecto a la concentración del pico.

\paragraph{Tanda de penaltis.}
Excluida del \textit{scope} causal por ser un deporte distinto. No
entra en \texttt{n\_shocks} ni en CATE; los intervalos y los
\textit{rankings} cubren regulación más prórroga.

\paragraph{Tarjetas, sustituciones y lesiones como \textit{shock}.}
Solo goles y \textit{near-misses} cuentan como tratamiento. Las
tarjetas, sustituciones y lesiones entran como \textit{flags} de
exclusión (\texttt{sub\_in\_window}, \texttt{truncated\_pre/post})
pero no como \textit{shock} propio.

\paragraph{Comportamiento colectivo del rival.}
El bloque \textit{leave-one-out} es del propio equipo del jugador.
La reacción del rival queda absorbida en el delta del bloque propio
y no se atribuye por separado.

\paragraph{Granularidad temporal de un minuto.}
Los moduladores continuos del contexto (minuto normalizado, leverage,
proximidad de eliminación, etc.) se evalúan a nivel de \textit{shock},
no a nivel de segundo. Cambios sub-minuto del contexto no entran en
el estimador.


\chapter{Validaciones empíricas adicionales}
\label{anexo:validaciones}

Este anexo recoge los diagnósticos de robustez que complementan los
reportados en el capítulo~\ref{ch:resultados}. Cada validación se
ejecuta sobre el panel jugador-\textit{shock}-minuto del Mundial
Qatar 2022 y se reporta en su tabla correspondiente.

\section{Test de placebo}
\label{anexo:val-placebo}

El test de placebo recoloca el \textit{shock} en instantes aleatorios
del partido distintos del gol real y vuelve a estimar el ATE por DiD
\textit{within-player} sobre cada celda canal-por-contexto. Si la
identificación causal es válida, el ATE real debe quedar fuera del
intervalo de confianza al 95\% de la distribución placebo
($H_0$: no hay efecto del \textit{shock}). La
tabla~\ref{tab:placebo} reporta el resultado.

\begin{table}[h]
  \centering
  \small
  \caption{Test de placebo sobre las ocho celdas
    canal-por-contexto.}
  \label{tab:placebo}
  \begin{tabular}{llrrrrl}
    \toprule
    Canal & Contexto & ATE real & Placebo media & Placebo SD & $p$ emp. & ATE $\in$ IC$_{95}$ \\
    \midrule
    Ofensivo  & Tras marcar  & $+0{,}00055$ & $+0{,}00004$ & $0{,}00136$ & $0{,}71$ & Sí \\
    Ofensivo  & Tras encajar & $+0{,}00000$ & $-0{,}00001$ & $0{,}00027$ & $0{,}98$ & Sí \\
    Defensivo & Tras marcar  & $-0{,}00023$ & $+0{,}00011$ & $0{,}00711$ & $0{,}96$ & Sí \\
    Defensivo & Tras encajar & $+0{,}00108$ & $-0{,}00014$ & $0{,}00635$ & $0{,}86$ & Sí \\
    Off-ball  & Tras marcar  & $-0{,}00002$ & $-0{,}00000$ & $0{,}00012$ & $0{,}91$ & Sí \\
    Off-ball  & Tras encajar & $+0{,}00001$ & $+0{,}00000$ & $0{,}00011$ & $0{,}96$ & Sí \\
    Físico    & Tras marcar  & $-0{,}00178$ & $+0{,}00106$ & $0{,}01179$ & $0{,}82$ & Sí \\
    Físico    & Tras encajar & $-0{,}01055$ & $-0{,}00011$ & $0{,}01266$ & $0{,}39$ & Sí \\
    \bottomrule
  \end{tabular}
\end{table}

El ATE real cae dentro del IC 95\% de la distribución placebo en las
ocho celdas, y el $p$ empírico a dos colas es superior a $0{,}39$ en
todas. La corrección de Benjamini-Hochberg sobre las ocho hipótesis
no produce ningún rechazo. La lectura es consistente con el ATE
poblacional indistinguible de cero reportado en el
capítulo~\ref{ch:resultados}: el test de placebo no rechaza la
nulidad porque el ATE real es ya cercano al ruido de la distribución
placebo, no porque la identificación esté mal especificada.

\section{Análisis de potencia ex-ante}
\label{anexo:val-power}

El análisis de potencia ex-ante calcula el \textit{Minimum Detectable
Effect} (MDE) al 80\% de potencia que admite el panel por celda,
ajustado por la correlación intra-cluster del estimador
\textit{cluster-robust} a nivel jugador (\textit{design effect}). La
tabla~\ref{tab:power} reporta el MDE efectivo por canal y contexto.

\begin{table}[h]
  \centering
  \small
  \caption{Potencia ex-ante por celda canal-por-contexto.}
  \label{tab:power}
  \begin{tabular}{llrrrr}
    \toprule
    Canal & Contexto & $n$ efectivo & ICC & MDE 80\% & ATE/MDE \\
    \midrule
    Ofensivo  & Tras marcar  & $469$ & $0{,}057$ & $0{,}00347$ & $0{,}16$ \\
    Ofensivo  & Tras encajar & $476$ & $0{,}000$ & $0{,}00074$ & $0{,}00$ \\
    Defensivo & Tras marcar  & $492$ & $0{,}000$ & $0{,}01904$ & $0{,}01$ \\
    Defensivo & Tras encajar & $433$ & $0{,}205$ & $0{,}01833$ & $0{,}06$ \\
    Off-ball  & Tras marcar  & $472$ & $0{,}048$ & $0{,}00031$ & $0{,}05$ \\
    Off-ball  & Tras encajar & $446$ & $0{,}140$ & $0{,}00030$ & $0{,}02$ \\
    Físico    & Tras marcar  & $492$ & $0{,}000$ & $0{,}03281$ & $0{,}05$ \\
    Físico    & Tras encajar & $460$ & $0{,}074$ & $0{,}03579$ & $0{,}29$ \\
    \bottomrule
  \end{tabular}
\end{table}

El MDE efectivo oscila entre $0{,}00030$ desviaciones estándar del
canal \textit{off-ball} y $0{,}03579$ del canal físico tras encajar.
El cociente ATE/MDE es inferior a $0{,}3$ en todas las celdas, lo
que sitúa los efectos observados muy por debajo del listón mínimo de
detección al 80\% de potencia con el tamaño de muestra disponible.
La potencia del panel para detectar efectos promedio en este corpus
es baja, y la lectura individual del CATE bayesiano recupera la
señal que el ATE poblacional no captura.

\section{Estratificación por fase del torneo}
\label{anexo:val-stage}

La estratificación por fase del torneo separa el ATE por grupos
(120 \textit{shocks}) y por fase eliminatoria (52 \textit{shocks}) y
verifica si la diferencia entre ambos contextos es sustantiva. La
tabla~\ref{tab:stage} reporta los coeficientes y los $p$-valores por
celda y por fase.

\begin{table}[h]
  \centering
  \small
  \caption{ATE estratificado por fase del torneo.}
  \label{tab:stage}
  \begin{tabular}{llrrrrrr}
    \toprule
    \multirow{2}{*}{Canal} & \multirow{2}{*}{Contexto} & \multicolumn{3}{c}{Fase de grupos} & \multicolumn{3}{c}{Fase eliminatoria} \\
    \cmidrule(lr){3-5} \cmidrule(lr){6-8}
                              &              & ATE & SE & $p$ & ATE & SE & $p$ \\
    \midrule
    Ofensivo  & Tras marcar  & $+0{,}00055$ & $0{,}00149$ & $0{,}71$ & $+0{,}00056$ & $0{,}00204$ & $0{,}78$ \\
    Ofensivo  & Tras encajar & $+0{,}00002$ & $0{,}00030$ & $0{,}94$ & $-0{,}00005$ & $0{,}00054$ & $0{,}93$ \\
    Defensivo & Tras marcar  & $-0{,}00086$ & $0{,}00734$ & $0{,}91$ & $+0{,}00125$ & $0{,}01487$ & $0{,}93$ \\
    Defensivo & Tras encajar & $+0{,}00109$ & $0{,}00785$ & $0{,}89$ & $+0{,}00106$ & $0{,}00995$ & $0{,}92$ \\
    Off-ball  & Tras marcar  & $+0{,}00001$ & $0{,}00015$ & $0{,}96$ & $-0{,}00007$ & $0{,}00014$ & $0{,}65$ \\
    Off-ball  & Tras encajar & $-0{,}00000$ & $0{,}00012$ & $0{,}99$ & $+0{,}00003$ & $0{,}00020$ & $0{,}88$ \\
    Físico    & Tras marcar  & $-0{,}00279$ & $0{,}01440$ & $0{,}85$ & $+0{,}00056$ & $0{,}02000$ & $0{,}98$ \\
    Físico    & Tras encajar & $-0{,}01456$ & $0{,}01518$ & $0{,}34$ & $-0{,}00121$ & $0{,}02237$ & $0{,}96$ \\
    \bottomrule
  \end{tabular}
\end{table}

Ningún coeficiente cruza el cinco por ciento de significatividad en
ninguna fase. La diferencia entre fase de grupos y fase eliminatoria
no es sustantiva en ninguna celda, lo que es consistente con la
identificación causal a nivel individual: la heterogeneidad
detectable del \CCV{} reside en el efecto individual del jugador
condicional al contexto (proximidad de eliminación, dirección del
\textit{shock}) y no en el promedio poblacional por fase del torneo.

\section{Diagnósticos de calibración del PSxG}
\label{anexo:val-psxg}

La tabla~\ref{tab:psxg-cal} reporta los diagnósticos completos de
calibración del modelo Post-Shot $xG$ sobre los dos conjuntos
relevantes: el \textit{out-of-fold} del entrenamiento StatsBomb
(3545 disparos, 136 partidos fuera del Mundial) y el \textit{holdout}
del Mundial Qatar 2022 (1494 disparos, 64 partidos). Se incluyen
las tres descomposiciones que la sección~\ref{sec:metricas-evaluacion}
del capítulo~\ref{ch:metodologia} formaliza: AUC, ECE binado, MCE,
Brier y su descomposición de Murphy en \textit{reliability},
\textit{resolution} y \textit{uncertainty}.

\begin{table}[h]
  \centering
  \small
  \caption{Diagnósticos de calibración del modelo PSxG.}
  \label{tab:psxg-cal}
  \begin{tabular}{llrrrrrr}
    \toprule
    Conjunto & Modelo & AUC & Brier & ECE & MCE & Rel. & Res. \\
    \midrule
    \textit{OOF}   & PSxG calibrado    & $0{,}968$ & $0{,}0415$ & $0{,}009$ & $0{,}67$ & $0{,}00009$ & $0{,}054$ \\
    \textit{OOF}   & PSxG \textit{raw} & $0{,}971$ & $0{,}0415$ & $0{,}013$ & $0{,}12$ & $0{,}00031$ & $0{,}055$ \\
    \textit{OOF}   & $xG$ StatsBomb    & $0{,}826$ & $0{,}0744$ & $0{,}009$ & $0{,}15$ & $0{,}00022$ & $0{,}020$ \\
    \midrule
    \textit{WC22}  & PSxG calibrado    & $0{,}976$ & $0{,}0372$ & $0{,}011$ & $0{,}13$ & $0{,}00018$ & $0{,}071$ \\
    \textit{WC22}  & PSxG \textit{raw} & $0{,}979$ & $0{,}0374$ & $0{,}012$ & $0{,}17$ & $0{,}00063$ & $0{,}072$ \\
    \textit{WC22}  & $xG$ StatsBomb    & $0{,}844$ & $0{,}0829$ & $0{,}017$ & $0{,}18$ & $0{,}00035$ & $0{,}030$ \\
    \bottomrule
  \end{tabular}
\end{table}

La AUC del PSxG calibrado sobre el \textit{holdout} del Mundial es
$0{,}976$, unos trece puntos por encima del $xG$ baseline. La
descomposición de Brier sitúa la \textit{reliability} dos órdenes de
magnitud por debajo de la \textit{resolution} sobre \textit{WC22}: la
desviación entre probabilidad predicha y frecuencia empírica por
\textit{bin} es despreciable frente a la capacidad discriminativa
agregada. La MCE sube de $0{,}12$ a $0{,}67$ tras calibrar, pero esa
métrica captura la peor desviación local sobre \textit{bins} pequeños
y se interpreta junto al ECE, que en cambio mejora de $0{,}013$ a
$0{,}009$. La transferencia al \textit{holdout} del Mundial preserva
la calibración (ECE $= 0{,}011$, \textit{reliability} $= 0{,}0002$),
las cifras que el capítulo~\ref{ch:resultados} reporta como ancla
técnica del modelo sobre el corpus objetivo.


\chapter{Nota técnica: espejo del tracking en prórroga}
\label{anexo:espejo-prorroga}

Este anexo documenta una inconsistencia detectada en una parte del
corpus PFF FC del Mundial Qatar 2022 entre el \textit{frame} de
tracking y el de \textit{eventing} durante los períodos de prórroga
(P3 y P4) y la corrección automática aplicada en el preprocesado.

\section{Naturaleza de la inconsistencia}
\label{anexo:espejo-naturaleza}

En una parte de los partidos del torneo, el \textit{frame} de
tracking entregado por PFF FC durante los períodos de prórroga
(período 3, primera parte de la prórroga; período 4, segunda parte)
está rotado 180\textdegree{} respecto al \textit{frame} del
\textit{eventing} del mismo instante. Es decir: el vector de
posiciones $(x, y)$ de los jugadores en el tracking aparece como
$(-x, -y)$ comparado con el sistema de coordenadas que el
\textit{eventing} utiliza para registrar la dirección de ataque de
cada equipo. La inconsistencia es consistente con la convención de
algunos proveedores de tracking de mantener la cámara orientada
hacia el mismo banderín durante la totalidad del partido, mientras
que el \textit{eventing} alterna la dirección de ataque por equipo
a cada cambio de período.

La inconsistencia no es universal en el corpus: afecta a una parte
de los partidos con prórroga pero no a todos. La verificación manual
en los partidos del torneo con prórroga ha identificado el patrón en
la final Argentina-Francia (\texttt{pff\_match\_id}~10517) y en
Japón-Croacia (\texttt{pff\_match\_id}~10506), y ha confirmado su
ausencia en Marruecos-España (\texttt{pff\_match\_id}~10508). La
detección caso a caso es la estrategia operativa adecuada.

\section{Mecanismo de detección automática}
\label{anexo:espejo-deteccion}

El preprocesado detecta la inconsistencia automáticamente comparando,
en el \textit{frame} de tracking inmediatamente posterior a la
transición de período, dos referencias independientes: la dirección
de ataque del equipo en posesión registrada en el \textit{eventing}
del primer evento del nuevo período, y la posición del portero del
equipo defensor en el \textit{frame} de tracking correspondiente. Si
las dos referencias son consistentes, el tracking se acepta tal cual;
si son opuestas, el \textit{pipeline} aplica una des-rotación de
180\textdegree{} ($x' = -x$, $y' = -y$) a todos los \textit{frames}
de los períodos afectados antes de cualquier cálculo posterior. La
detección es silenciosa y produce un \textit{flag} \textit{per}-partido
que documenta la corrección aplicada.

\section{Ámbito de aplicación de la corrección}
\label{anexo:espejo-ambito}

La corrección se aplica en todas las capas del \textit{pipeline} que
mezclan coordenadas del tracking con la dirección de ataque del
equipo: la estimación del OBSO y C-OBSO en el canal \textit{off-ball},
la atribución al defensor más cercano al balón en el canal
defensivo y la superficie de \textit{pitch control} sobre la que se
ilustra el caso de Mbappé en la figura~\ref{fig:ppcf-mbappe} del
capítulo~\ref{ch:desarrollo}. Las capas que utilizan métricas
invariantes a rotación ---el cálculo de velocidades, distancias y
\textit{sprints} del canal físico, y los descriptores escalares de la
identificación de \textit{shocks}--- no se ven afectadas por la
inconsistencia y operan sobre el tracking crudo sin des-rotación.

La sincronización entre el \textit{frame} de tracking y el evento
correspondiente se realiza mediante el campo \texttt{start\_frame}
del \textit{eventing}, no mediante el reloj de juego
\texttt{periodGameClockTime}, que en algunos partidos del corpus
deriva por unas décimas de segundo durante la prórroga y produciría
desalineamientos sistemáticos en los cálculos \textit{frame-level}.


\chapter{Uso de asistencia de IA generativa en la elaboración del trabajo}
\label{anexo:uso-ia}

En cumplimiento de las buenas prácticas académicas sobre transparencia
en el uso de herramientas de IA generativa, este anexo declara
explícitamente el alcance con el que estas herramientas se han
utilizado durante la elaboración del presente trabajo.

\section{Alcance del uso}
\label{anexo:uso-ia-alcance}

La asistencia ha provenido del modelo Claude Sonnet 4.6 de Anthropic,
empleado de forma acotada en dos roles operativos durante la
elaboración del trabajo. El primero es la asistencia en aspectos
formales de la redacción de la memoria ---sintaxis castellana, claridad
expositiva sobre formulaciones matemáticas, consistencia terminológica
entre capítulos, formato \LaTeX{} y depuración de errores de
compilación. El segundo es el apoyo al desarrollo del \textit{pipeline}
de código del proyecto ---revisión de implementaciones contra los
\textit{papers} de referencia, generación de comentarios explicativos
sobre bloques de cálculo, identificación de fallos lógicos y propuesta
de \textit{snippets} de validación empírica.

\section{Lo que no ha hecho la IA}
\label{anexo:uso-ia-noha}

Se declara explícitamente que la asistencia de IA generativa no ha
intervenido en los siguientes aspectos sustantivos del trabajo. La
formulación de la pregunta de investigación, la definición conceptual
del marco \CCV, el diseño de la estrategia de identificación causal
(las cinco capas formalizadas en el capítulo~\ref{ch:metodologia}) y
la selección de los \textit{building blocks} del estado del arte son
contribución intelectual exclusiva del autor. La narrativa de cada
capítulo, la estructura argumental, la interpretación de los
resultados y las conclusiones sustantivas son de redacción original
del autor; la asistencia en sintaxis no ha llegado en ningún caso a
ser redacción autónoma de párrafos completos sobre ideas no
articuladas previamente por el autor.

\section{Validación de los outputs asistidos}
\label{anexo:uso-ia-validacion}

Todos los outputs sobre los que la IA ha asistido han sido verificados
contra fuentes primarias antes de incorporarse al trabajo. Las cifras
numéricas reportadas en el cuerpo del trabajo se han validado
ejecutando el código sobre los datos reales del Mundial Qatar 2022 y
contrastando contra los \textit{parquet} de \texttt{outputs/} y
\texttt{data/parquet/derived/}; las citas bibliográficas se han
verificado individualmente contra el DOI o el repositorio original;
las fórmulas matemáticas se han contrastado contra los \textit{papers}
de referencia citados en el capítulo~\ref{ch:estado-arte}; las
implementaciones de código que la IA ha contribuido a revisar se han
validado mediante \textit{tests} de consistencia interna y mediante
reproducción de cifras agregadas publicadas por los autores de los
\textit{papers} originales. La responsabilidad última sobre la
corrección de cada afirmación del trabajo es exclusivamente del autor.
